\section*{Modelling with Algebra}
\subsection*{Application of Algebra}

\subsubsection*{Introduction}

This section is a continuation of what was discussed in Year 1, sets and binomial
expansion. De Morgan's Laws is an important theorem in sets as its application
transcends beyond union, intersection and complements. These laws are
central to Boolean Algebra where it is used to simplify expressions. In Digital
Circuit Design, it is used to transform logic gates with the intention of not only
optimising digital circuitry but also reducing cost and saving space. In Computer
programming, it is used to simplify logical expressions in code, which can make
code cleaner, simpler, and more readable. Database systems use these laws to
optimise queries in SQL (Structured Query Language), which enhances data
retrieval more efficiently.

The Binomial Theorem is a strong tool that makes the process of expanding
binomial expressions much simpler. In developing this theorem, you further
your algebra skills and prepare yourselves for higher classes of mathematics.
The Binomial Theorem is applied in finance, investment, probability, statistics,
computer science, genetics, engineering, game theory, marketing, and physics. It
helps in calculating compound interest, risk assessment, modelling success/failure
outcomes, understanding data structure complexities, predicting inheritance
probabilities and enhancing decision-making in real-world scenarios.

\begin{keyideas}
\begin{itemize}
  \item \textbf{Binomial Coefficients:} The coefficients $^{n}C_{r}$ represent the number of ways
  to choose $r$ elements from $n$ elements and can be found using Pascal's
  Triangle.
  \[
    {}^{n}C_{r} = \frac{n!}{(n-r)!\,r!}
  \]
  \item \textbf{Terms in the Expansion:} Each term in the expansion is of the form:
  \[
    {}^{n}C_{r}\, a^{\,n-r} b^{\,r}
  \]
  \item The Binomial Theorem provides a formula for expanding expressions of
  the form $(a+b)^n$ where $n$ is a non-negative integer.
  \item The complement of an empty set ($\phi$) is equal to the universal set ($\mu$), i.e.\ $\phi' = \mu$.
  \item The intersection of a set $A$ and its complement ($A'$) is equal to the null set
  ($\varnothing$), i.e.\ $A \cap A' = \varnothing$.
  \item The union of a set $A$ and its complement ($A'$) is equal to the universal set
  ($\mu$), i.e.\ $A \cup A' = \mu$.
\end{itemize}
\end{keyideas}
